\chapter*{Notar}
\addcontentsline{toc}{chapter}{Notar}
\markboth{Notar}{}

{\footnotesize

\noteentry{Om tilhøvet til traktaten.}{Denne boka er den destruktive halvdelen av eit par. Traktaten \textsc{Formlære} legg fram fundamentet i full aksiomatisk form, med proposisjonar, empirisk fundament og falsifiseringsvilkår; \textsc{Grunnlaust} river ned påstanden om at faget alt har eit fundament. Påstandane i kapittel \textsc{ix} er difor ikkje grunngjevne her; dei er grunngjevne der. Den som vil felle rammeverket, skal felle det i traktaten, der vilkåra for å felle det står lista.}

\noteentry{Om diagnose framfor skuld.}{Argumentet gjennom heile boka er strukturelt, ikkje moralsk. Påstanden er aldri at ein einskild lærar, skule eller designar har gjort noko gale. Påstanden er at disiplinen manglar eit fundament, at dette er eit kollektivt fråvær og ikkje ei individuell forsøming, og at konsekvensane difor berre kan rettast kollektivt og strukturelt. Kvar setning som kan lesast som ein påstand om at namngjevne personar har feila, er meint strukturelt og bør lesast slik.}

\noteentry{Om den tomme formelen (kap. \textsc{ii}).}{Argumentet følgjer Michl\kref{1,\,19} sin disjunksjon: «funksjon» tyder anten noko snevert, og då er forma underbestemt, eller alt, og då er påstanden tom. Den logiske strukturen er den til ein tautologi. Det positive alternativet — at forma er ein rest etter mange samtidige trykk — er traktaten sitt, og generaliserer Michls negative innsikt til eit positivt rammeverk.}

\noteentry{Om vrange problem (kap. \textsc{v}).}{Rittel og Webber\kref{4} og Buchanan\kref{5} har rett om at mange designproblem er vrange. Innvendinga gjeld bruken: at vrangskapen vert nytta som grunn til at teori er umogleg, når han eigentleg er spesifikasjonen på kva teorien må handtere. Biologien sitt tilpassingslandskap\kref{9,\,12} er det historiske motdømet — eit presist apparat for nett mange samtidige, motstridande trykk.}

\noteentry{Om avleggar-samanlikninga (kap. \textsc{viii}).}{Samanlikninga med frenologi, alkemi og vitalisme er strukturell, ikkje retorisk: det desse felta delte med designfaget er katalog utan teori, mønstergjenkjenning utan fundament. Kuhn\kref{15} sitt omgrep om det før-paradigmatiske faget gjev den presise diagnosen. Det nye er at agentisk materiale\kref{14} og generative system no fjernar den tredje vegen — å halde fram som før.}

\bigskip

}

\chapter*{Referansar}
\addcontentsline{toc}{chapter}{Referansar}
\markboth{Referansar}{}

\refentry{[1] Michl, J. (1995). Form follows WHAT? The modernist notion of function as a carte blanche. \textit{1. Designforum Finland Magazine}, 1.}
\refentry{[2] Schön, D. A. (1983). \textit{The Reflective Practitioner: How Professionals Think in Action}. Basic Books.}
\refentry{[3] Simon, H. A. (1969). \textit{The Sciences of the Artificial}. MIT Press.}
\refentry{[4] Rittel, H. W. J. \& Webber, M. M. (1973). Dilemmas in a general theory of planning. \textit{Policy Sciences}, 4(2), 155--169.}
\refentry{[5] Buchanan, R. (1992). Wicked problems in design thinking. \textit{Design Issues}, 8(2), 5--21.}
\refentry{[6] Cross, N. (1982). Designerly ways of knowing. \textit{Design Studies}, 3(4), 221--227.}
\refentry{[7] Frayling, C. (1993). Research in art and design. \textit{Royal College of Art Research Papers}, 1(1), 1--5.}
\refentry{[8] Stiny, G. (2006). \textit{Shape: Talking about Seeing and Doing}. MIT Press.}
\refentry{[9] Stiny, G. \& Gips, J. (1972). Shape grammars and the generative specification of painting and sculpture. \textit{Information Processing 71}, 1460--1465.}
\refentry{[10] Hatchuel, A. \& Weil, B. (2009). C-K design theory: An advanced formulation. \textit{Research in Engineering Design}, 19(4), 181--192.}
\refentry{[11] Thompson, D'A. W. (1917). \textit{On Growth and Form}. Cambridge University Press.}
\refentry{[12] Wright, S. (1932). The roles of mutation, inbreeding, crossbreeding, and selection in evolution. \textit{Proc. Sixth Int. Congress of Genetics}, 1, 356--366.}
\refentry{[13] Fields, C. \& Levin, M. (2022). Competency in navigating arbitrary spaces as an invariant for analyzing cognition in diverse systems. \textit{Entropy}, 24(6), 819.}
\refentry{[14] Levin, M. (2022). Technological approach to mind everywhere: An experimentally-grounded framework for understanding diverse bodies and minds. \textit{Frontiers in Systems Neuroscience}, 16, 768201.}
\refentry{[15] Kuhn, T. S. (1962). \textit{The Structure of Scientific Revolutions}. University of Chicago Press.}
\refentry{[16] Gropius, W. (1919). \textit{Manifest und Programm des Staatlichen Bauhauses}. Staatliches Bauhaus Weimar.}
\refentry{[17] Norman, D. A. (2013). \textit{The Design of Everyday Things} (Revised ed.). Basic Books.}
\refentry{[18] Krippendorff, K. (2006). \textit{The Semantic Turn: A New Foundation for Design}. CRC Press.}
\refentry{[19] Michl, J. (2014). A case against the modernist regime in design education. \textit{Archnet-IJAR: International Journal of Architectural Research}, 8(2), 36--46.}
\refentry{[20] Biddle, S. (2019--2024). Palantir and ICE reporting series. \textit{The Intercept}. Sjå òg Mijente, \textit{Who's Behind ICE?}}
\refentry{[21] Abraham, Y. (2024). «Lavender»: AI-baserte målrettingssystem i Gaza. \textit{+972 Magazine / Local Call}.}
\refentry{[22] Haidt, J. (2024). \textit{The Anxious Generation: How the Great Rewiring of Childhood Is Causing an Epidemic of Mental Illness}. Penguin Press.}
